% Options for packages loaded elsewhere
\PassOptionsToPackage{unicode}{hyperref}
\PassOptionsToPackage{hyphens}{url}
\documentclass[
]{article}
\usepackage{xcolor}
\usepackage{amsmath,amssymb}
\setcounter{secnumdepth}{-\maxdimen} % remove section numbering
\usepackage{iftex}
\ifPDFTeX
  \usepackage[T1]{fontenc}
  \usepackage[utf8]{inputenc}
  \usepackage{textcomp} % provide euro and other symbols
\else % if luatex or xetex
  \usepackage{unicode-math} % this also loads fontspec
  \defaultfontfeatures{Scale=MatchLowercase}
  \defaultfontfeatures[\rmfamily]{Ligatures=TeX,Scale=1}
\fi
\usepackage{lmodern}
\ifPDFTeX\else
  % xetex/luatex font selection
\fi
% Use upquote if available, for straight quotes in verbatim environments
\IfFileExists{upquote.sty}{\usepackage{upquote}}{}
\IfFileExists{microtype.sty}{% use microtype if available
  \usepackage[]{microtype}
  \UseMicrotypeSet[protrusion]{basicmath} % disable protrusion for tt fonts
}{}
\makeatletter
\@ifundefined{KOMAClassName}{% if non-KOMA class
  \IfFileExists{parskip.sty}{%
    \usepackage{parskip}
  }{% else
    \setlength{\parindent}{0pt}
    \setlength{\parskip}{6pt plus 2pt minus 1pt}}
}{% if KOMA class
  \KOMAoptions{parskip=half}}
\makeatother
\usepackage{color}
\usepackage{fancyvrb}
\newcommand{\VerbBar}{|}
\newcommand{\VERB}{\Verb[commandchars=\\\{\}]}
\DefineVerbatimEnvironment{Highlighting}{Verbatim}{commandchars=\\\{\}}
% Add ',fontsize=\small' for more characters per line
\newenvironment{Shaded}{}{}
\newcommand{\AlertTok}[1]{\textcolor[rgb]{1.00,0.00,0.00}{\textbf{#1}}}
\newcommand{\AnnotationTok}[1]{\textcolor[rgb]{0.38,0.63,0.69}{\textbf{\textit{#1}}}}
\newcommand{\AttributeTok}[1]{\textcolor[rgb]{0.49,0.56,0.16}{#1}}
\newcommand{\BaseNTok}[1]{\textcolor[rgb]{0.25,0.63,0.44}{#1}}
\newcommand{\BuiltInTok}[1]{\textcolor[rgb]{0.00,0.50,0.00}{#1}}
\newcommand{\CharTok}[1]{\textcolor[rgb]{0.25,0.44,0.63}{#1}}
\newcommand{\CommentTok}[1]{\textcolor[rgb]{0.38,0.63,0.69}{\textit{#1}}}
\newcommand{\CommentVarTok}[1]{\textcolor[rgb]{0.38,0.63,0.69}{\textbf{\textit{#1}}}}
\newcommand{\ConstantTok}[1]{\textcolor[rgb]{0.53,0.00,0.00}{#1}}
\newcommand{\ControlFlowTok}[1]{\textcolor[rgb]{0.00,0.44,0.13}{\textbf{#1}}}
\newcommand{\DataTypeTok}[1]{\textcolor[rgb]{0.56,0.13,0.00}{#1}}
\newcommand{\DecValTok}[1]{\textcolor[rgb]{0.25,0.63,0.44}{#1}}
\newcommand{\DocumentationTok}[1]{\textcolor[rgb]{0.73,0.13,0.13}{\textit{#1}}}
\newcommand{\ErrorTok}[1]{\textcolor[rgb]{1.00,0.00,0.00}{\textbf{#1}}}
\newcommand{\ExtensionTok}[1]{#1}
\newcommand{\FloatTok}[1]{\textcolor[rgb]{0.25,0.63,0.44}{#1}}
\newcommand{\FunctionTok}[1]{\textcolor[rgb]{0.02,0.16,0.49}{#1}}
\newcommand{\ImportTok}[1]{\textcolor[rgb]{0.00,0.50,0.00}{\textbf{#1}}}
\newcommand{\InformationTok}[1]{\textcolor[rgb]{0.38,0.63,0.69}{\textbf{\textit{#1}}}}
\newcommand{\KeywordTok}[1]{\textcolor[rgb]{0.00,0.44,0.13}{\textbf{#1}}}
\newcommand{\NormalTok}[1]{#1}
\newcommand{\OperatorTok}[1]{\textcolor[rgb]{0.40,0.40,0.40}{#1}}
\newcommand{\OtherTok}[1]{\textcolor[rgb]{0.00,0.44,0.13}{#1}}
\newcommand{\PreprocessorTok}[1]{\textcolor[rgb]{0.74,0.48,0.00}{#1}}
\newcommand{\RegionMarkerTok}[1]{#1}
\newcommand{\SpecialCharTok}[1]{\textcolor[rgb]{0.25,0.44,0.63}{#1}}
\newcommand{\SpecialStringTok}[1]{\textcolor[rgb]{0.73,0.40,0.53}{#1}}
\newcommand{\StringTok}[1]{\textcolor[rgb]{0.25,0.44,0.63}{#1}}
\newcommand{\VariableTok}[1]{\textcolor[rgb]{0.10,0.09,0.49}{#1}}
\newcommand{\VerbatimStringTok}[1]{\textcolor[rgb]{0.25,0.44,0.63}{#1}}
\newcommand{\WarningTok}[1]{\textcolor[rgb]{0.38,0.63,0.69}{\textbf{\textit{#1}}}}
\usepackage{longtable,booktabs,array}
\newcounter{none} % for unnumbered tables
\usepackage{calc} % for calculating minipage widths
% Correct order of tables after \paragraph or \subparagraph
\usepackage{etoolbox}
\makeatletter
\patchcmd\longtable{\par}{\if@noskipsec\mbox{}\fi\par}{}{}
\makeatother
% Allow footnotes in longtable head/foot
\IfFileExists{footnotehyper.sty}{\usepackage{footnotehyper}}{\usepackage{footnote}}
\makesavenoteenv{longtable}
\setlength{\emergencystretch}{3em} % prevent overfull lines
\providecommand{\tightlist}{%
  \setlength{\itemsep}{0pt}\setlength{\parskip}{0pt}}
\usepackage{bookmark}
\IfFileExists{xurl.sty}{\usepackage{xurl}}{} % add URL line breaks if available
\urlstyle{same}
\hypersetup{
  hidelinks,
  pdfcreator={LaTeX via pandoc}}

\author{}
\date{}

\begin{document}

\section{Cryptographic Operations Audit
Report}\label{cryptographic-operations-audit-report}

\textbf{Repository:} \texttt{/Users/z/work/lux/standard/contracts/}
\textbf{Date:} 2026-01-30 \textbf{Auditor:} Automated Security Analysis

\begin{center}\rule{0.5\linewidth}{0.5pt}\end{center}

\subsection{Executive Summary}\label{executive-summary}

This audit examines all cryptographic operations across the Lux Standard
contracts. The codebase implements a diverse set of cryptographic
schemes including ECDSA, FROST Schnorr signatures, Lamport signatures,
EIP-712 typed data, and Merkle proofs. While most implementations follow
best practices, several findings require attention.

\textbf{Severity Summary:}

\begin{itemize}
\tightlist
\item
  CRITICAL: 2
\item
  HIGH: 4
\item
  MEDIUM: 6
\item
  LOW: 8
\item
  INFORMATIONAL: 5
\end{itemize}

\begin{center}\rule{0.5\linewidth}{0.5pt}\end{center}

\subsection{1. Signature Verification
(ecrecover)}\label{1-signature-verification-ecrecover}

\subsubsection{1.1 Files Using
ecrecover}\label{11-files-using-ecrecover}

{\def\LTcaptype{none} % do not increment counter
\begin{longtable}[]{@{}lll@{}}
\toprule\noalign{}
File & Line & Context \\
\midrule\noalign{}
\endhead
\bottomrule\noalign{}
\endlastfoot
\texttt{bridge/Teleport.sol} & 273 & MPC oracle signature
verification \\
\texttt{bridge/Bridge.sol} & 327 & MPC oracle signature verification \\
\texttt{governance/DAO.sol} & 323 & Vote by signature \\
\texttt{safe/FROST.sol} & 15-26 & Schnorr signature via ecrecover
trick \\
\end{longtable}
}

\subsubsection{1.2 Findings}\label{12-findings}

\paragraph{CRITICAL: Missing Signature Malleability Protection in Bridge
Contracts}\label{critical-missing-signature-malleability-protection-in-bridge-contracts}

\textbf{Files:} \texttt{bridge/Teleport.sol}, \texttt{bridge/Bridge.sol}

\textbf{Issue:} Both contracts use raw \texttt{ecrecover} without
checking the \texttt{s} value is in the lower half of the curve order.
ECDSA signatures are malleable: for any valid signature
\texttt{(r,\ s,\ v)}, the signature
\texttt{(r,\ n\ -\ s,\ v\textquotesingle{})} is also valid for the same
message.

\begin{Shaded}
\begin{Highlighting}[]
\NormalTok{// bridge/Teleport.sol:271{-}273}
\NormalTok{function recoverSigner(bytes32 message, bytes memory sig) internal pure returns (address) \{}
\NormalTok{    (uint8 v, bytes32 r, bytes32 s) = splitSignature(sig);}
\NormalTok{    return ecrecover(message, v, r, s);  // No malleability check}
\NormalTok{\}}
\end{Highlighting}
\end{Shaded}

\textbf{Impact:} An attacker could submit a malleable signature variant
to bypass replay protection if the signature bytes themselves are used
as the replay key (which they are in
\texttt{transactionMap{[}signedTXInfo{]}}).

\textbf{Recommendation:} Use OpenZeppelin\textquotesingle s
\texttt{ECDSA.recover()} which enforces
\texttt{s\ \textless{}=\ secp256k1n/2}, or add explicit check:

\begin{Shaded}
\begin{Highlighting}[]
\NormalTok{require(uint256(s) \textless{}= 0x7FFFFFFFFFFFFFFFFFFFFFFFFFFFFFFF5D576E7357A4501DDFE92F46681B20A0);}
\end{Highlighting}
\end{Shaded}

\paragraph{HIGH: No Zero Address Check After
ecrecover}\label{high-no-zero-address-check-after-ecrecover}

\textbf{Files:} \texttt{bridge/Teleport.sol}, \texttt{bridge/Bridge.sol}

\textbf{Issue:} \texttt{ecrecover} returns \texttt{address(0)} on
invalid signatures rather than reverting. The bridge contracts do check
the recovered address against \texttt{MPCOracleAddrMap}, but this relies
on \texttt{address(0)} not being in the oracle map. An explicit check
should be added.

\begin{Shaded}
\begin{Highlighting}[]
\NormalTok{// bridge/Bridge.sol:445{-}451}
\NormalTok{address signer = recoverSigner(}
\NormalTok{    prefixed(keccak256(abi.encodePacked(message))),}
\NormalTok{    signedTXInfo\_}
\NormalTok{);}
\NormalTok{// Missing: require(signer != address(0), "Invalid signature");}
\NormalTok{require(MPCOracleAddrMap[signer].exists, "Unauthorized Signature");}
\end{Highlighting}
\end{Shaded}

\textbf{Recommendation:} Add explicit zero address check immediately
after recovery.

\paragraph{LOW: DAO.sol Uses Raw ecrecover
(Acceptable)}\label{low-daosol-uses-raw-ecrecover-acceptable}

\textbf{File:} \texttt{governance/DAO.sol:323}

The DAO contract uses \texttt{ecrecover} directly but the context is
acceptable because:

\begin{enumerate}
\def\labelenumi{\arabic{enumi}.}
\tightlist
\item
  The signature is used to identify the voter, not for replay protection
\item
  Invalid signatures (returning \texttt{address(0)}) are caught by
  subsequent \texttt{require(signer\ !=\ address(0))}
\end{enumerate}

\begin{center}\rule{0.5\linewidth}{0.5pt}\end{center}

\subsection{2. Hash Functions
(keccak256)}\label{2-hash-functions-keccak256}

\subsubsection{2.1 Hash Collision Risks}\label{21-hash-collision-risks}

\paragraph{MEDIUM: abi.encodePacked with Multiple Dynamic
Types}\label{medium-abiencodepacked-with-multiple-dynamic-types}

\textbf{Multiple Files}

Using \texttt{abi.encodePacked} with multiple dynamic-length arguments
can lead to hash collisions:

\begin{Shaded}
\begin{Highlighting}[]
\NormalTok{// bridge/Teleport.sol:183}
\NormalTok{return string(abi.encodePacked(amt, toTargetAddrStr, txid, tokenAddrStrHash, chainIdStr, vault));}
\end{Highlighting}
\end{Shaded}

\textbf{Example Collision:}

\begin{itemize}
\tightlist
\item
  \texttt{encodePacked("a",\ "bc")} = \texttt{encodePacked("ab",\ "c")}
  = \texttt{"abc"}
\end{itemize}

\textbf{Affected Locations:}

\begin{itemize}
\tightlist
\item
  \texttt{bridge/Teleport.sol:183} - Message construction
\item
  \texttt{bridge/Bridge.sol:276-284} - Message construction
\item
  \texttt{privacy/BulletproofVerifier.sol:166-167} - Challenge
  computation
\item
  \texttt{account/Account.sol:206} - Session key message hash
\end{itemize}

\textbf{Recommendation:} Use \texttt{abi.encode} instead of
\texttt{abi.encodePacked} for all security-critical hashing, or ensure
at least one fixed-length element separates dynamic elements.

\paragraph{LOW: Proper Usage Patterns
Observed}\label{low-proper-usage-patterns-observed}

Several contracts correctly use \texttt{abi.encode} for EIP-712
structured data:

\begin{itemize}
\tightlist
\item
  \texttt{governance/DAO.sol:305-320} - Domain separator and struct hash
\item
  \texttt{governance/Council.sol:224-243} - Transaction hashing
\item
  \texttt{tokens/LRC20/extensions/LRC20Permit.sol:60-61} - Permit hash
\item
  \texttt{router/IntentRouter.sol:444-459} - Order hash
\end{itemize}

\begin{center}\rule{0.5\linewidth}{0.5pt}\end{center}

\subsection{3. Nonce Management}\label{3-nonce-management}

\subsubsection{3.1 Files with Nonce
Management}\label{31-files-with-nonce-management}

{\def\LTcaptype{none} % do not increment counter
\begin{longtable}[]{@{}lll@{}}
\toprule\noalign{}
File & Nonce Type & Implementation \\
\midrule\noalign{}
\endhead
\bottomrule\noalign{}
\endlastfoot
\texttt{account/Account.sol} & Session nonce & Incremented after use \\
\texttt{tokens/LRC20/extensions/LRC20Permit.sol} & EIP-2612 nonce &
OpenZeppelin Nonces \\
\texttt{fhe/governance/ConfidentialLRC20Votes.sol} & Delegation nonce &
Custom mapping \\
\texttt{router/IntentRouter.sol} & Order nonce & Invalidation mapping \\
\texttt{safe/SafeThresholdLamportModule.sol} & Lamport nonce &
Incremented after use \\
\texttt{bridge/teleport/TeleportBridge.sol} & Burn nonce & Sequential
counter \\
\end{longtable}
}

\subsubsection{3.2 Findings}\label{32-findings}

\paragraph{HIGH: Nonce Not Validated Before Use in
Account.sol}\label{high-nonce-not-validated-before-use-in-accountsol}

\textbf{File:} \texttt{account/Account.sol:206}

\begin{Shaded}
\begin{Highlighting}[]
\NormalTok{function executeWithSession(...) external returns (bytes memory) \{}
\NormalTok{    bytes32 hash = keccak256(abi.encodePacked(target, value, data, nonce));}
\NormalTok{    // ... signature verification ...}
\NormalTok{    nonce++;  // Nonce incremented AFTER use}
\NormalTok{\}}
\end{Highlighting}
\end{Shaded}

\textbf{Issue:} The nonce is included in the hash but
there\textquotesingle s no validation that the signature was created
with the current nonce. An attacker could potentially replay old
signatures if the same \texttt{(target,\ value,\ data)} tuple recurs.

\textbf{Recommendation:} Include explicit nonce validation or use a
nonce mapping per signer.

\paragraph{MEDIUM: Potential Nonce Desync in
ConfidentialLRC20Votes}\label{medium-potential-nonce-desync-in-confidentiallrc20votes}

\textbf{File:}
\texttt{fhe/governance/ConfidentialLRC20Votes.sol:137-139}

\begin{Shaded}
\begin{Highlighting}[]
\NormalTok{if (nonce != nonces[delegator]++) \{}
\NormalTok{    revert SignatureNonceInvalid();}
\NormalTok{\}}
\end{Highlighting}
\end{Shaded}

\textbf{Issue:} The nonce is incremented even when the signature
verification fails (due to the \texttt{++} operator in the condition).
If \texttt{SignatureChecker.isValidSignatureNow} fails first (line 134),
the nonce is not affected, but the order of checks could lead to
inconsistent state if reordered.

\textbf{Recommendation:} Separate nonce validation from nonce increment
for clarity.

\paragraph{LOW: Good Practice - TeleportBridge Nonce
Management}\label{low-good-practice---teleportbridge-nonce-management}

\textbf{File:} \texttt{bridge/teleport/TeleportBridge.sol:156}

\begin{Shaded}
\begin{Highlighting}[]
\NormalTok{uint256 currentNonce = burnNonce++;}
\end{Highlighting}
\end{Shaded}

Good pattern: sequential nonce incremented atomically, included in
canonical burn ID computation.

\begin{center}\rule{0.5\linewidth}{0.5pt}\end{center}

\subsection{4. Randomness Sources}\label{4-randomness-sources}

\subsubsection{4.1 Analysis of block.timestamp/block.number
Usage}\label{41-analysis-of-blocktimestampblocknumber-usage}

\paragraph{MEDIUM: block.timestamp Used in Pool ID
Generation}\label{medium-blocktimestamp-used-in-pool-id-generation}

\textbf{File:} \texttt{mocks/MockZChainAMM.sol:36}

\begin{Shaded}
\begin{Highlighting}[]
\NormalTok{poolId = keccak256(abi.encodePacked(assetA, assetB, block.timestamp));}
\end{Highlighting}
\end{Shaded}

\textbf{Issue:} Using \texttt{block.timestamp} for generating
identifiers allows miners/validators to influence the outcome. This is a
mock contract, but the pattern should not be used in production.

\paragraph{MEDIUM: Potential Time-Based Manipulation in Oracle
Staleness}\label{medium-potential-time-based-manipulation-in-oracle-staleness}

\textbf{File:} \texttt{precompile/interfaces/IOracle.sol:299}

\begin{Shaded}
\begin{Highlighting}[]
\NormalTok{require(block.timestamp {-} p.timestamp \textless{}= maxAge, "Oracle: stale price");}
\end{Highlighting}
\end{Shaded}

\textbf{Issue:} Price staleness checks rely on \texttt{block.timestamp}
which can be manipulated by validators within bounds (\textasciitilde15
seconds on Ethereum). For high-value transactions, this tolerance may be
exploitable.

\textbf{Recommendation:} Consider using multiple oracle sources or
implementing TWAP for price-sensitive operations.

\begin{center}\rule{0.5\linewidth}{0.5pt}\end{center}

\subsection{5. EIP-712 Typed Data
Signing}\label{5-eip-712-typed-data-signing}

\subsubsection{5.1 Implementation
Review}\label{51-implementation-review}

{\def\LTcaptype{none} % do not increment counter
\begin{longtable}[]{@{}lll@{}}
\toprule\noalign{}
Contract & Domain Separator & Findings \\
\midrule\noalign{}
\endhead
\bottomrule\noalign{}
\endlastfoot
\texttt{governance/DAO.sol} & Dynamic construction & Computed per-call
(gas inefficient but correct) \\
\texttt{governance/Council.sol} & DOMAIN\_SEPARATOR\_TYPEHASH &
Correctly typed \\
\texttt{tokens/LRC20/extensions/LRC20Permit.sol} & OpenZeppelin EIP712 &
Best practice implementation \\
\texttt{router/IntentRouter.sol} & OpenZeppelin EIP712 & Best practice
implementation \\
\texttt{bridge/teleport/TeleportBridge.sol} & OpenZeppelin EIP712 & Best
practice implementation \\
\texttt{fhe/access/PermissionedV2.sol} & Modified EIP712 & Intentionally
omits verifyingContract \\
\texttt{account/Account.sol} & Defined but unused & DOMAIN\_TYPEHASH
defined but EIP-712 not fully implemented \\
\end{longtable}
}

\subsubsection{5.2 Findings}\label{52-findings}

\paragraph{HIGH: DAO.sol Domain Separator Lacks
Version}\label{high-daosol-domain-separator-lacks-version}

\textbf{File:} \texttt{governance/DAO.sol:305-311}

\begin{Shaded}
\begin{Highlighting}[]
\NormalTok{bytes32 domainSeparator = keccak256(}
\NormalTok{    abi.encode(}
\NormalTok{        keccak256("EIP712Domain(string name,uint256 chainId,address verifyingContract)"),}
\NormalTok{        keccak256(bytes("LuxDAO")),}
\NormalTok{        block.chainid,}
\NormalTok{        address(this)}
\NormalTok{    )}
\NormalTok{);}
\end{Highlighting}
\end{Shaded}

\textbf{Issue:} The domain separator is missing the \texttt{version}
field that is standard in EIP-712. While functional, this deviates from
the standard and may cause compatibility issues with signing libraries.

\textbf{Recommendation:} Add version field:

\begin{Shaded}
\begin{Highlighting}[]
\NormalTok{keccak256("EIP712Domain(string name,string version,uint256 chainId,address verifyingContract)")}
\end{Highlighting}
\end{Shaded}

\paragraph{LOW: Domain Separator Not Cached in
DAO.sol}\label{low-domain-separator-not-cached-in-daosol}

\textbf{File:} \texttt{governance/DAO.sol}

The domain separator is recomputed on every \texttt{castVoteBySig} call.
While correct (handles chain forks), it\textquotesingle s gas
inefficient. Consider caching with fork detection like
OpenZeppelin\textquotesingle s implementation.

\paragraph{INFORMATIONAL: PermissionedV2 Intentionally Omits
verifyingContract}\label{informational-permissionedv2-intentionally-omits-verifyingcontract}

\textbf{File:} \texttt{fhe/access/PermissionedV2.sol}

This is intentional design to allow single signatures to work across
multiple contracts. The access control is handled via the
\texttt{contracts} and \texttt{projects} arrays in the permission
struct. This is documented and acceptable.

\begin{center}\rule{0.5\linewidth}{0.5pt}\end{center}

\subsection{6. Permit Signatures
(ERC-2612)}\label{6-permit-signatures-erc-2612}

\subsubsection{6.1 Implementations}\label{61-implementations}

{\def\LTcaptype{none} % do not increment counter
\begin{longtable}[]{@{}ll@{}}
\toprule\noalign{}
Contract & Base Implementation \\
\midrule\noalign{}
\endhead
\bottomrule\noalign{}
\endlastfoot
\texttt{tokens/LRC20/extensions/LRC20Permit.sol} & Custom (follows
standard) \\
\texttt{governance/Stake.sol} & OpenZeppelin ERC20Permit \\
\texttt{liquid/LiquidLUX.sol} & OpenZeppelin ERC20Permit \\
\texttt{tokens/LRC4626/LRC4626.sol} & OpenZeppelin ERC20Permit \\
\texttt{tokens/LRC20/LRC20.sol} & OpenZeppelin ERC20Permit \\
\end{longtable}
}

\subsubsection{6.2 Findings}\label{62-findings}

\paragraph{LOW: LRC20Permit Uses ECDSA.recover
Correctly}\label{low-lrc20permit-uses-ecdsarecover-correctly}

\textbf{File:} \texttt{tokens/LRC20/extensions/LRC20Permit.sol:65}

\begin{Shaded}
\begin{Highlighting}[]
\NormalTok{address signer = ECDSA.recover(hash, v, r, s);}
\end{Highlighting}
\end{Shaded}

Good: Uses OpenZeppelin\textquotesingle s ECDSA library which includes
malleability protection.

\paragraph{INFORMATIONAL: All Permit Implementations Follow
ERC-2612}\label{informational-all-permit-implementations-follow-erc-2612}

All permit implementations correctly:

\begin{itemize}
\tightlist
\item
  Use sequential nonces via \texttt{\_useNonce()}
\item
  Include deadline validation
\item
  Use EIP-712 typed data hashing
\item
  Validate signer matches owner
\end{itemize}

\begin{center}\rule{0.5\linewidth}{0.5pt}\end{center}

\subsection{7. Merkle Proof
Verification}\label{7-merkle-proof-verification}

\subsubsection{7.1 Implementations
Found}\label{71-implementations-found}

{\def\LTcaptype{none} % do not increment counter
\begin{longtable}[]{@{}ll@{}}
\toprule\noalign{}
Contract & Purpose \\
\midrule\noalign{}
\endhead
\bottomrule\noalign{}
\endlastfoot
\texttt{privacy/ZNote.sol} & Note commitment tree \\
\texttt{privacy/ShieldedTreasury.sol} & Commitment tree \\
\texttt{privacy/Poseidon2Commitments.sol} & ZK commitment tree \\
\texttt{privacy/PrivateBridge.sol} & Cross-chain commitment \\
\texttt{precompile/interfaces/IZK.sol} & Precompile interface \\
\texttt{precompile/interfaces/IHash.sol} & Hash precompile interface \\
\end{longtable}
}

\subsubsection{7.2 Findings}\label{72-findings}

\paragraph{MEDIUM: Second Preimage Attack Risk in Custom Merkle
Implementations}\label{medium-second-preimage-attack-risk-in-custom-merkle-implementations}

\textbf{Files:} \texttt{privacy/ZNote.sol},
\texttt{privacy/ShieldedTreasury.sol}

\begin{Shaded}
\begin{Highlighting}[]
\NormalTok{// privacy/ZNote.sol:382{-}385}
\NormalTok{if (pathIndices[i] == 0) \{}
\NormalTok{    currentHash = keccak256(abi.encodePacked(currentHash, proof[i]));}
\NormalTok{\} else \{}
\NormalTok{    currentHash = keccak256(abi.encodePacked(proof[i], currentHash));}
\NormalTok{\}}
\end{Highlighting}
\end{Shaded}

\textbf{Issue:} Without domain separation between leaf and internal
nodes, second preimage attacks may be possible. An attacker could
potentially construct a proof where a leaf hash equals an internal node
hash.

\textbf{Recommendation:} Use different hash domains for leaves vs
internal nodes:

\begin{Shaded}
\begin{Highlighting}[]
\NormalTok{// For leaves:}
\NormalTok{keccak256(abi.encodePacked(bytes1(0x00), leafData))}
\NormalTok{// For internal nodes:}
\NormalTok{keccak256(abi.encodePacked(bytes1(0x01), left, right))}
\end{Highlighting}
\end{Shaded}

\paragraph{INFORMATIONAL: OpenZeppelin MerkleProof Not
Used}\label{informational-openzeppelin-merkleproof-not-used}

The contracts implement custom Merkle verification rather than using
OpenZeppelin\textquotesingle s audited \texttt{MerkleProof} library.
Custom implementations are acceptable for specific use cases (e.g.,
incremental Merkle trees) but increase audit surface.

\begin{center}\rule{0.5\linewidth}{0.5pt}\end{center}

\subsection{8. Commitment Schemes}\label{8-commitment-schemes}

\subsubsection{8.1 Implementations}\label{81-implementations}

{\def\LTcaptype{none} % do not increment counter
\begin{longtable}[]{@{}ll@{}}
\toprule\noalign{}
Contract & Scheme Type \\
\midrule\noalign{}
\endhead
\bottomrule\noalign{}
\endlastfoot
\texttt{safe/SafeLSSSigner.sol} & LSS commitment \\
\texttt{privacy/Poseidon2Commitments.sol} & Pedersen-style \\
\texttt{privacy/PrivateTeleport.sol} & Note commitments \\
\texttt{precompile/interfaces/IZK.sol} & KZG, FRI commitments \\
\end{longtable}
}

\subsubsection{8.2 Findings}\label{82-findings}

\paragraph{LOW: Commitment Binding in
SafeLSSSigner}\label{low-commitment-binding-in-safelsssigner}

\textbf{File:} \texttt{safe/SafeLSSSigner.sol:134}

\begin{Shaded}
\begin{Highlighting}[]
\NormalTok{commitment: keccak256(abi.encodePacked(\_threshold, \_totalSigners, \_publicKey))}
\end{Highlighting}
\end{Shaded}

The commitment binds threshold parameters to public key. This is a good
pattern for preventing parameter manipulation.

\begin{center}\rule{0.5\linewidth}{0.5pt}\end{center}

\subsection{9. FROST Schnorr
Signatures}\label{9-frost-schnorr-signatures}

\subsubsection{9.1 Implementation
Analysis}\label{91-implementation-analysis}

\textbf{File:} \texttt{safe/FROST.sol}

The FROST library implements RFC 9591 FROST(secp256k1, SHA-256) Schnorr
signatures using the \texttt{ecrecover} precompile trick for mul-mul-add
operations.

\subsubsection{9.2 Findings}\label{92-findings}

\paragraph{INFORMATIONAL: Signature Malleability
Handled}\label{informational-signature-malleability-handled}

\textbf{File:} \texttt{safe/FROST.sol:294-311}

\begin{Shaded}
\begin{Highlighting}[]
\NormalTok{// TODO(nlordell): I don\textquotesingle{}t think this is required for Schnorr}
\NormalTok{// signatures, but do it anyway just in case.}
\NormalTok{\{}
\NormalTok{    bool pOk = isValidPublicKey(px, py);}
\NormalTok{    bool rOk = \_isOnCurve(rx, ry);}
\NormalTok{    bool zOk = \_isScalar(z);}
\NormalTok{    // ...}
\NormalTok{\}}
\end{Highlighting}
\end{Shaded}

Good: The implementation validates that \texttt{z} is a valid scalar (in
range \texttt{(0,\ N)}), preventing trivial malleability via \texttt{z}
manipulation.

\paragraph{LOW: Public Key Restriction
Documented}\label{low-public-key-restriction-documented}

\textbf{File:} `safe/FROST.sol:228-230**

\begin{Shaded}
\begin{Highlighting}[]
\NormalTok{/// @dev Note that public key\textquotesingle{}s x{-}coordinate \textasciigrave{}px\textasciigrave{} must be smaller than the}
\NormalTok{/// curve order for the math trick with \textasciigrave{}ecrecover\textasciigrave{} to work.}
\end{Highlighting}
\end{Shaded}

The restriction that \texttt{px\ \textless{}\ N} is documented and
enforced in \texttt{isValidPublicKey()}. This is a known limitation of
the ecrecover trick approach.

\begin{center}\rule{0.5\linewidth}{0.5pt}\end{center}

\subsection{10. Lamport Signatures}\label{10-lamport-signatures}

\subsubsection{10.1 Implementation
Analysis}\label{101-implementation-analysis}

\textbf{File:} \texttt{safe/SafeThresholdLamportModule.sol}

Implements one-time Lamport signatures for post-quantum security with
M-Chain MPC threshold control (LP-134; the legacy ``T-Chain MPC'' name
referred to today's M-Chain bridge-custody chain).

\subsubsection{10.2 Findings}\label{102-findings}

\paragraph{INFORMATIONAL: Proper One-Time Key
Rotation}\label{informational-proper-one-time-key-rotation}

\textbf{File:} \texttt{safe/SafeThresholdLamportModule.sol:197-202}

\begin{Shaded}
\begin{Highlighting}[]
\NormalTok{bytes32 oldPkh = pkh;}
\NormalTok{pkh = nextPKH;}
\NormalTok{lamportNonce++;}

\NormalTok{emit LamportKeyRotated(oldPkh, nextPKH);}
\end{Highlighting}
\end{Shaded}

Good: Key rotation is enforced as part of the execution flow, ensuring
one-time property of Lamport signatures.

\paragraph{LOW: Domain Separation Properly
Implemented}\label{low-domain-separation-properly-implemented}

\textbf{File:} `safe/SafeThresholdLamportModule.sol:180-185**

\begin{Shaded}
\begin{Highlighting}[]
\NormalTok{uint256 m = uint256(keccak256(abi.encodePacked(}
\NormalTok{    safeTxHash,}
\NormalTok{    nextPKH,}
\NormalTok{    address(this),   // Prevent cross{-}contract replay}
\NormalTok{    block.chainid    // Prevent cross{-}chain replay}
\NormalTok{)));}
\end{Highlighting}
\end{Shaded}

Good: Domain separation includes contract address and chain ID to
prevent replay attacks across chains and contracts.

\begin{center}\rule{0.5\linewidth}{0.5pt}\end{center}

\subsection{11. Replay Attack Vectors}\label{11-replay-attack-vectors}

\subsubsection{11.1 Analysis}\label{111-analysis}

\paragraph{CRITICAL: Signature-Based Replay Protection in Legacy
Bridges}\label{critical-signature-based-replay-protection-in-legacy-bridges}

\textbf{Files:} \texttt{bridge/Teleport.sol}, \texttt{bridge/Bridge.sol}

\begin{Shaded}
\begin{Highlighting}[]
\NormalTok{// bridge/Bridge.sol:441{-}443}
\NormalTok{require(}
\NormalTok{    !transactionMap[signedTXInfo\_].exists,}
\NormalTok{    "Duplicated Transaction Hash"}
\NormalTok{);}
\end{Highlighting}
\end{Shaded}

\textbf{Issue:} Replay protection uses the raw signature bytes as the
key. Due to ECDSA signature malleability (see Section 1.2), an attacker
can derive an alternative valid signature for the same message and
bypass replay protection.

\textbf{Attack Scenario:}

\begin{enumerate}
\def\labelenumi{\arabic{enumi}.}
\tightlist
\item
  User submits valid bridge mint with signature \texttt{(r,\ s,\ v)}
\item
  Attacker computes malleable signature
  \texttt{(r,\ n-s,\ v\textquotesingle{})} for same message
\item
  Attacker replays with malleable signature, passing replay check
\item
  Tokens minted twice
\end{enumerate}

\textbf{Recommendation:}

\begin{enumerate}
\def\labelenumi{\arabic{enumi}.}
\tightlist
\item
  Use message hash (not signature) as replay key, OR
\item
  Enforce \texttt{s} is in lower half of curve order before storing
\end{enumerate}

\paragraph{LOW: Proper Replay Protection
Examples}\label{low-proper-replay-protection-examples}

\begin{itemize}
\tightlist
\item
  \texttt{bridge/teleport/TeleportBridge.sol}: Uses \texttt{claimId}
  computed from claim data, not signature
\item
  \texttt{router/IntentRouter.sol}: Uses
  \texttt{filledAmounts{[}orderHash{]}} where orderHash is from order
  data
\item
  \texttt{governance/DAO.sol}: Uses \texttt{receipt.hasVoted} mapping
  keyed by voter address
\end{itemize}

\begin{center}\rule{0.5\linewidth}{0.5pt}\end{center}

\subsection{12. Front-Running
Considerations}\label{12-front-running-considerations}

\subsubsection{12.1 Analysis}\label{121-analysis}

\paragraph{MEDIUM: Bridge Mints Susceptible to
Front-Running}\label{medium-bridge-mints-susceptible-to-front-running}

\textbf{Files:} \texttt{bridge/Teleport.sol}, \texttt{bridge/Bridge.sol}

Once an MPC oracle signature is observed in the mempool, anyone can
submit the mint transaction. The recipient is fixed in the signed
message, but gas price wars could occur.

\textbf{Mitigation:} The contracts use \texttt{nonReentrant} which
prevents reentrancy but not front-running. Consider implementing
commit-reveal or private mempool submission for high-value transfers.

\paragraph{LOW: Intent Router Has MEV Protection
Documentation}\label{low-intent-router-has-mev-protection-documentation}

\textbf{File:} \texttt{router/IntentRouter.sol}

The contract explicitly mentions "MEV protection via private mempools"
in its documentation, indicating awareness of front-running concerns.

\begin{center}\rule{0.5\linewidth}{0.5pt}\end{center}

\subsection{13. Recommendations
Summary}\label{13-recommendations-summary}

\subsubsection{Critical Priority}\label{critical-priority}

\begin{enumerate}
\def\labelenumi{\arabic{enumi}.}
\tightlist
\item
  \textbf{Fix signature malleability in bridge contracts} - Add
  \texttt{s} value validation or use OpenZeppelin ECDSA
\item
  \textbf{Change replay protection key from signature to message hash}
  in legacy bridges
\end{enumerate}

\subsubsection{High Priority}\label{high-priority}

\begin{enumerate}
\def\labelenumi{\arabic{enumi}.}
\setcounter{enumi}{2}
\tightlist
\item
  \textbf{Add explicit \texttt{address(0)} check after ecrecover} in all
  signature verification
\item
  \textbf{Add version field to DAO.sol EIP-712 domain separator}
\item
  \textbf{Fix nonce validation in Account.sol} session key execution
\item
  \textbf{Review ConfidentialLRC20Votes nonce increment ordering}
\end{enumerate}

\subsubsection{Medium Priority}\label{medium-priority}

\begin{enumerate}
\def\labelenumi{\arabic{enumi}.}
\setcounter{enumi}{6}
\tightlist
\item
  \textbf{Replace \texttt{abi.encodePacked} with \texttt{abi.encode}}
  for security-critical hashing
\item
  \textbf{Add leaf/node domain separation} in custom Merkle
  implementations
\item
  \textbf{Review oracle staleness timing} for manipulation resistance
\item
  \textbf{Avoid block.timestamp in identifier generation} (even in
  mocks, as bad patterns propagate)
\end{enumerate}

\subsubsection{Low Priority}\label{low-priority}

\begin{enumerate}
\def\labelenumi{\arabic{enumi}.}
\setcounter{enumi}{10}
\tightlist
\item
  Consider caching domain separator in DAO.sol
\item
  Document public key restrictions for FROST signatures in user-facing
  docs
\item
  Add comments explaining nonce patterns in complex flows
\end{enumerate}

\begin{center}\rule{0.5\linewidth}{0.5pt}\end{center}

\subsection{14. Positive Findings}\label{14-positive-findings}

The following patterns demonstrate good cryptographic practices:

\begin{enumerate}
\def\labelenumi{\arabic{enumi}.}
\tightlist
\item
  \textbf{OpenZeppelin Usage}: Most newer contracts use audited
  OpenZeppelin libraries (ECDSA, EIP712, Nonces)
\item
  \textbf{EIP-712 Adoption}: Widespread use of typed data signing for
  user experience and security
\item
  \textbf{Domain Separation}: Most signature schemes include chain ID
  and contract address
\item
  \textbf{Sequential Nonces}: Proper sequential nonce patterns in permit
  and order systems
\item
  \textbf{One-Time Key Rotation}: Lamport implementation correctly
  enforces key rotation
\item
  \textbf{Snapshot Voting}: DAO uses checkpointed voting to prevent
  flash loan attacks
\end{enumerate}

\begin{center}\rule{0.5\linewidth}{0.5pt}\end{center}

\subsection{15. Files Reviewed}\label{15-files-reviewed}

\subsubsection{Primary Cryptographic
Contracts}\label{primary-cryptographic-contracts}

\begin{itemize}
\tightlist
\item
  \texttt{bridge/Teleport.sol}
\item
  \texttt{bridge/Bridge.sol}
\item
  \texttt{bridge/teleport/TeleportBridge.sol}
\item
  \texttt{governance/DAO.sol}
\item
  \texttt{governance/Council.sol}
\item
  \texttt{account/Account.sol}
\item
  \texttt{safe/FROST.sol}
\item
  \texttt{safe/SafeThresholdLamportModule.sol}
\item
  \texttt{tokens/LRC20/extensions/LRC20Permit.sol}
\item
  \texttt{router/IntentRouter.sol}
\item
  \texttt{fhe/governance/ConfidentialLRC20Votes.sol}
\item
  \texttt{fhe/access/PermissionedV2.sol}
\end{itemize}

\subsubsection{Supporting Files}\label{supporting-files}

\begin{itemize}
\tightlist
\item
  \texttt{privacy/ZNote.sol}
\item
  \texttt{privacy/ShieldedTreasury.sol}
\item
  \texttt{privacy/Poseidon2Commitments.sol}
\item
  \texttt{precompile/interfaces/IZK.sol}
\item
  \texttt{precompile/interfaces/IHash.sol}
\item
  Various yield strategy and bridge token contracts
\end{itemize}

\begin{center}\rule{0.5\linewidth}{0.5pt}\end{center}

\textbf{End of Report}

\end{document}
